\section{Limitations and future work}
\label{app:limits}

All quantitative conclusions concern CLIP vision encoders. The headline uses $k=8$, $m_C=64$,
and the $2\sigma$ resolution convention; we sweep $k$, $m_C$, and the resolution multiplier below.
ViT-L/14 covers only \texttt{q} and \texttt{v}, and $C_H$ is estimated from seven of the twenty
benchmark datasets under one recorded mixture. One deliberately different composition is tested
below; broader calibration compositions and language models remain open.
The reported finite-pool stability intervals hold the calibration draw, fold partition, block
rule, and $m_C$ fixed; they do not incorporate uncertainty from estimating or selecting the null.
They are not superpopulation coverage statements: each task contributes one checkpoint/training
realisation, so training-run and checkpoint variation are unmeasured.
Scalar activation-null values use Eq.~\ref{eq:null-expectation} exactly; adapter intervals hold
their pooled $64$-draw realised Monte Carlo null values fixed. We draw $50{,}000$ crossed
task--layer resampling replicates. Independent multinomial endpoint counts $C_t$ and layer counts $L_\ell$
give each row weight $C_aC_bL_\ell$; all module types within a selected layer remain grouped.
Each replicate takes the ratio after weighted averaging. The pseudorandom seed is $20260915$.
As a finite-pool influence check for the exact activation arms, leave-one-task-out reproduced
fractions range $24.685$--$25.896\%$, $21.856$--$22.939\%$, and $11.310$--$12.165\%$ for
ViT-B/16, ViT-B/32, and ViT-L/14; the corresponding leave-one-layer-out ranges are
$24.359$--$28.411\%$, $21.272$--$26.392\%$, and $10.955$--$13.587\%$. Every residual remains positive.

The raw outer-product gradient identity involves activations along each fine-tuning trajectory
and, under scalar-step gradient descent, constrains the update to their span. Our null instead
conditions on the frozen base's second moment. Feature drift, adaptive coordinatewise
preconditioning, and weight decay break an exact passage between those statements. Consequently,
the activation-conditioned decomposition is descriptive: its post-fine-tuning occupancy profile
can include learned structure, and identifying its origin would require an intervention or an
additional stability assumption that we do not test.

The LoRA result is limited to rank $16$, one shared-start protocol, eight tasks, and the reported
attention modules. Its factor null preserves $A$ and the singular spectrum of $B$, not the
product spectrum of $BA$. In a four-module, six-task stress test the attainable plant delivers
$37.4$--$57.0\%$ of its nominal increment, and the null recovers $100.7$--$101.1\%$ of that
delivered signal. This establishes level and response only for that planting operator; the
$87.5\%$ projector-basis fraction has no dedicated positive control. A controlled
shared-versus-independent-$A_0$ experiment is absent, and $B$ remains about five times above
isotropic chance. Finally, calibration is nearly rank-preserving and the cross-term basis does
not improve compressed-merge accuracy; the evidence supports measurement and attribution claims,
not a better merging algorithm.

\section{Invariants of the overlap statistic}
\label{app:invariants}

Every numerical property below is checked directly in the analysis; these checks should accompany
the released measurement code rather than be inferred from library defaults.

Beyond the scalar occupancies and energies in \S\ref{sec:null}, the group action fixes
$\Delta W_iP_b\Delta W_i^\top$ and the eigenvalues of
$V_i^{(k)\top}P_bV_i^{(k)}$ for every block. The former contains output-side block Gram geometry
and the latter its principal-angle spectrum. A looser null conditioning only on $t_{ib}$ could
therefore give a different attribution.

\paragraph{Why Eq.~\ref{eq:null-expectation} is exact.}
Haar averaging within block $b$ gives
$\mathbb E[Q_i^\top\Pi_iQ_i]=\sum_b(t_{ib}/d_b)P_b$, where
$d_b=\operatorname{rank}(P_b)$. For independent task-wise draws, linearity and
$P_bP_c=\mathbf 1\{b=c\}P_b$ give
$\mathbb E\operatorname{tr}(Q_i^\top\Pi_iQ_iQ_j^\top\Pi_jQ_j)
=\sum_b t_{ib}t_{jb}/d_b$; division by $k$ yields Eq.~\ref{eq:null-expectation}.

\begin{itemize}
\item \textbf{Orientation.} \texttt{top\_right} returns the right singular vectors, of shape
$(d_{\mathrm{in}}, k)$, matching an explicit SVD on a rectangular matrix where a transpose would
be silent. The convention $\Delta W = U S V^\top$ puts $V$ on the input side; swapping it
reverses every conclusion in this paper.
\item \textbf{Normalisation.} $\mathcal{O}(A,A)=1$ to $10^{-6}$, and two independent uniformly
random $k$-subspaces average $0.03995$ against $k/d = 0.04000$ over $300$ draws. We confirm the
same identity at $(d,k) \in \{(200,8),(200,32),(400,8),(3584,8),(100,10)\}$.
\item \textbf{Degenerate inputs are refused.} The SVD of a zero matrix returns an arbitrary
orthonormal basis, and returns the \emph{same} one each time, so two zero task vectors score
$\mathcal{O}=1.0000$. In a LoRA checkpoint every module the adapter does not target is exactly
zero, so an unguarded run over an adapter pool fills its averages with the maximum possible
value taken from modules that did not change. \texttt{top\_right} raises on such an input and
the caller drops and counts those cells.
\item \textbf{Checkpoint reads.} Weights are read by matching module paths to checkpoint keys by
suffix, since \texttt{CLIPVisionModel.named\_modules()} yields \texttt{encoder.layers.0...}\ while
the checkpoint stores \texttt{vision\_model.encoder.layers.0...}. A suffix matching more than one key is refused: the full CLIP checkpoint holds both
\texttt{text\_model.*} and \texttt{vision\_model.*} under identical suffixes, and picking either
silently would build task vectors from the text tower and compare them against a vision
covariance. Tensors returned this way are bit-identical to a \texttt{state\_dict} load.
\end{itemize}

\section{Calibration choices for whitening and activation blocks}
\label{app:calibration}

\paragraph{Block-error rule.}
Partition the calibration set into eight disjoint folds and compute the ordered covariance
eigenvalues $\lambda_j^{(f)}$ independently on each fold. With their mean $\bar\lambda_j^{F}$,
the implementation uses the split-sample standard error
$\sigma_j=[\sum_{f=1}^{8}(\lambda_j^{(f)}-\bar\lambda_j^{F})^2/(8\cdot7)]^{1/2}$.
The blocks are the transitive consecutive components induced by overlap of adjacent
$\widehat\lambda_j\mathbin{\pm}2\sigma_j$ intervals, where $\widehat\lambda_j$ is the eigenvalue
from the full calibration pool. This is an operational coarsening, not a confidence interval for
eigenspace identifiability.

\paragraph{Finite calibration refits.}
Holding the seven source counts and extraction procedure fixed, two source-stratified bootstrap
refits (seeds $20260916$ and $20260917$) give full ViT-B/16 exact reproduced fractions
$25.102\%$ and $25.302\%$, versus the $25.345\%$ headline, for a maximum $0.243$-percentage-point
shift. This is the observed range of two refits, not a calibration confidence interval, and does
not test a different source composition.

\paragraph{Composition stress control.}
Separately, we sample without replacement from deterministic test splits to balance $334$ MNIST,
$334$ Stanford Cars, and $333$ SVHN images (seed $20260919$). This deliberately different
calibration gives exact ViT-B/16 null $0.029621$, residual $0.074953$, and reproduced fraction
$21.482\%$, or $3.863$ percentage points below the seven-source headline. Holding this covariance
fixed, the crossed finite-pool stability intervals are $[16.8975,26.2690]\%$ for the fraction and
$[+0.048606,+0.109015]$ for the residual. The median participation ratio is $12.691$, with
$71/72$ modules below $64$, so both composition and spectral effective rank have changed. This is
a stress control, not an estimate of calibration uncertainty or an interval over possible compositions.

A complementary post-hoc stratification asks whether zero, one, or both endpoints of a task pair
belong to the seven calibration sources ($78/91/21$ pairs). The reproduced fractions rise from
$23.06/26.56/28.33\%$ for ViT-B/16, $20.38/23.45/24.55\%$ for ViT-B/32, and
$10.20/12.64/14.47\%$ for ViT-L/14; the residual remains positive in every cell. This gradient is
consistent with composition sensitivity, but it does not isolate calibration membership from task
identity and is not used for inference.

\paragraph{Truncation.} Whitening two task vectors by a shared $C_H$ manufactures overlap when
it is aggressive. On independent task vectors whose true overlap is chance ($0.0400$), a
Tikhonov ridge at $\lambda = 10^{-4}\bar{w}$ returns $0.7199$ and RegMean's published
off-diagonal shrinkage $\alpha = 0.9$ returns $0.0931$. The mechanism is that $C_H^{-1/2}$
amplifies small eigendirections and drives both task vectors onto the same few of them: at
$\lambda = 10^{-4}\bar{w}$, $96.5\%$ of the whitened read mass lands on the twenty smallest
stored eigendirections. Whitening only the top $m_{\rm white}$ directions keeps reference overlap
near chance ($0.0394$, $0.0405$, $0.0422$ at $m_{\rm white}=8,32,64$).

We set $m_{\rm white}$ per module from that module's participation ratio
$(\sum_i w_i)^2 / \sum_i w_i^2$, used here as an operational effective-rank truncation. In our
calibration refits, directions beyond this cutoff were unstable; the participation ratio is not
an identifiability threshold. The ratio must be computed from the full spectrum:
taken over only the $m_C$ retained eigenvalues it is bounded above by
$m_C$ by Cauchy--Schwarz, which understates it and makes any gate comparing
the two a tautology.

\paragraph{Why a block count was rejected.} \S\ref{sec:method} sets the rotation blocks from
what the calibration data can resolve. We record here what the alternative looked like, because
the alternative is what a reader would reach for first.

Fixing a block count and complement rotation rank leaves two choices that no invariant selects,
and the answer moves with them: sweeping the rotation rank of the complement from $32$
to $512$ moved the reproduced fraction from $48\%$ to $23\%$ under the null as it then stood.
Subtracting that null's offset on a constructed reference did not stabilise it. The
resolution-based construction eliminates those two outcome-sensitive choices, since the
complement is a single unresolved block and is rotated in full. It does not eliminate design
choices: the uncertainty multiplier, fold-SE scheme, calibration sample, and stored dimension
remain part of the null specification.

The sweep is also where the defect of \S\ref{sec:defects} hid in plain sight. A rotation rank
that changes the answer by a factor of two is a rotation rank that is destroying different
amounts of structure, and the endpoint the construction actually calls for, rotating the whole
complement, lies outside the range that was swept.

\paragraph{Block structure is a function of the calibration budget.}
Blocks are the groups of eigendirections that our operational interval rule treats as unresolved.
They tend to sharpen as the calibration set grows: estimated eigenvalue errors contract, so fewer
adjacent intervals overlap and blocks split. This behavior does not make interval overlap a
confidence theorem for eigenspace identifiability.
The first three columns below are measured on nested subsets of one image pool. At each budget we
partition images into eight disjoint folds and compute the eight foldwise eigenvalue replicates
used by the split-sample rule in \S\ref{sec:method}. The remainder scales the measured standard errors,
the pool having stopped at $1001$.

\begin{table}[h]
\centering
\begin{tabular}{@{}lrrrcrrr@{}}
\toprule
& \multicolumn{3}{c}{measured} & & \multicolumn{3}{c}{extrapolated} \\
\cmidrule(r){2-4}\cmidrule(l){6-8}
calibration images & $224$ & $500$ & $1001$ & & $2000$ & $5000$ & $\infty$ \\
\midrule
ViT-B/16 & $16$ & $25$ & $37$ & & $46$ & $55$ & $64$ \\
ViT-B/32 & $15$ & $21$ & $37$ & & $44$ & $55$ & $64$ \\
\bottomrule
\end{tabular}
\end{table}

This is a resource-resolution trade-off that the method makes visible. Too few images and the blocks are coarse, the
null rotates across gaps that are in fact resolved, and its constructed-reference offset grows.
As the budget grows, the rule tends to split blocks. Under a simple population spectrum and fixed
$m_C=64$, its retained-region limit is $\{\pm1\}^{64}$ while the unresolved complement remains
$O(d_{\rm in}-64)$. Thus the limit is not the identity: signs randomise off-diagonal terms within
the retained region, and the complement is fully rotated. It has no power against agreement
expressed only as shared retained-axis occupancy, while retaining power against sign-sensitive or
complement orientation alignment. The $\infty$ column extrapolates the retained block count only;
it is not a measured power result.

The budget matters for a reason that is visible in the block counts themselves. A run at $224$
images drawn independently of the pool above resolved $13$ and $11$ blocks where the nested $224$
subset resolves $16$ and $15$, so at that budget the block structure is not stable across draws
of the calibration set, let alone across architectures. At $1001$ the two architectures happen
to agree at a median of $37$ blocks and the same median largest block. This is a useful
cross-architecture check, not proof of sampling stability. The calibration budget is part of the
measurement specification and must be fixed and reported.

The reproduced fraction is not reported against the budget, since the two budgets were measured
under nulls that differ by more than the budget: the $224$-image run predates the corrections
of \S\ref{sec:defects}, so comparing them would charge the calibration set for a change that
belongs to the null.

\paragraph{Where the blocked region ends.}
The activation-null stored dimension $m_C$ divides the space the null treats two ways: inside
$\mathrm{span}(V)$ it rotates within resolvable blocks, outside it rotates in full. We sweep
$k\in\{4,8,16,32\}$, $m_C\in\{16,32,64\}$, and uncertainty multiplier
$z\in\{1,2,3\}$ on each reported pool:

\begin{table}[h]
\centering
\begin{tabular}{@{}lrr@{}}
\toprule
pool & $k=8$, all $(m_C,z)$ & full $(k,m_C,z)$ grid \\
\midrule
ViT-B/16 & $20.22$--$25.61\%$ & $13.28$--$30.37\%$ \\
ViT-B/32 & $17.94$--$22.55\%$ & $13.31$--$28.51\%$ \\
ViT-L/14 $q/v$ & $5.84$--$12.00\%$ & $5.36$--$19.01\%$ \\
\bottomrule
\end{tabular}
\end{table}

At fixed $(k,m_C)$, the largest observed change from varying $z$ is $0.754$ percentage points; most of
the fixed-$k$ range comes from $m_C$. This is variation in the conditional estimand---storing more
directions changes what the null holds fixed---rather than numerical instability of its exact
estimator. We pre-specify the headline at $m_C=64,z=2$; $m_C$ remains distinct from the per-module
whitening truncation $m_{\rm white}$ above.

\paragraph{The calibration must use the observed spectrum.} The reference offset is not a property
of the null alone; it depends strongly on the spectrum of the task vectors the null is
calibrated against. Substituting a synthetic exponential spectrum for the observed one raises
the realised reference offset at eight blocks from $+0.0041$ to $+0.0254$, a factor of six.
This sensitivity is another reason not to transfer that diagnostic into the estimand.
The reason is that real task
vectors have rapidly decaying spectra, which makes the top-$k$ subspace well determined and the
block rotation effective, whereas a flat spectrum leaves that subspace ill-determined and the
block structure dominates. A null calibrated against a spectrum the data does not have
calibrates the wrong quantity.

\section{LoRA factor audit}
\label{app:lora-factors}

\begin{table}[h]
\centering
\caption{ViT-B/16, eight tasks, $28$ pairs, all $24$ targeted attention modules; all entries use top-$8$ singular
subspaces. $B$ is initialised at zero and is therefore entirely learned; $A$ is not.
\texttt{lin-LoRA} is the partial linearization of \citet{tang2024linearization}, which leaves
$A$ bit-identical across tasks.}
\begin{tabular}{@{}lrrr@{}}
\toprule
& $A$ right & $B$ left & $\Delta W$ right \\
\midrule
LoRA-16   & $0.8759$ & $0.0540$ & $0.6075$ \\
lin-LoRA  & $1.0000$ & $0.0570$ & $0.6649$ \\
isotropic & $0.0104$ & $0.0104$ & $0.0104$ \\
\bottomrule
\end{tabular}
\end{table}

At the $k=8$ boundary, the relative singular gap
$(s_8-s_9)/s_8$ is above $10^{-3}$ for every observed $A$ and $BA$ cell in both
eight-task, $24$-module adapter pools.  Across $24{,}576$ factor-null $BA$ redraws, only four
fall below $10^{-3}$ and none below $10^{-4}$; the median gap is about $0.099$ in each arm.
Thus exact or numerical cutoff ties do not drive the aggregate factor-null result, although the
finite-$k$ choice remains part of the estimand.

\section{Power of the LoRA null}
\label{app:lora-power}

\begin{table}[h]
\centering
\caption{A rank-$16$ LoRA stress test, ViT-B/16, four attention modules, six tasks. The fitted
factor null has pre-planting overlap $\mathcal O_0=0.4825$. From an unconstrained plant $P_s$ we
form the attainable $P_s^\star=P_sA^+A=(\alpha/r)B_sA$, where
$B_s=(r/\alpha)P_sA^+$. Both raw and null terms score $P_s^\star$. Delivery divides
$\Delta\mathrm{raw}$ by $(s/k)(1-\mathcal O_0)$; recovery is level-centred as in
Eq.~\ref{eq:recovery}.}
\label{tab:lora-power}
\begin{tabular}{@{}lrrrrr@{}}
\toprule
planted & raw & delivery & initial $X_s$ & corrected $X_s$ & recovery \\
\midrule
$0/8$ & $0.4825$ & ---        & $-0.1485$ & $-0.0073$ & --- \\
$1/8$ & $0.5067$ & $37.386\%$ & $-0.1243$ & $+0.0172$ & $101.117\%$ \\
$2/8$ & $0.5469$ & $49.790\%$ & $-0.0840$ & $+0.0579$ & $101.078\%$ \\
$4/8$ & $0.6301$ & $57.043\%$ & $-0.0009$ & $+0.1414$ & $100.701\%$ \\
\bottomrule
\end{tabular}
\end{table}

The initial norm-matched generator also has approximately $100\%$ level-centred response, but a
large $s=0$ error. Spectrum matching changes that level from $-0.1485$ to $-0.0073$ while
retaining near-full response. The correction therefore fixes level; it is not evidence that the
initial null had low power or upward bias. Values just above $100\%$ reflect finite draws.

\section{Merged accuracy by rank budget}
\label{app:merge}

\begin{table}[h]
\centering
\caption{Merged accuracy averaged over the eight tasks, ViT-B/32, $500$ test images per task,
update compressed to rank $r$. Zero-shot is $38.9\%$ and uncompressed task arithmetic is
$59.5\%$.}
\label{tab:merge}
\begin{tabular}{@{}lrrrrrr@{}}
\toprule
rank $r$ & $8$ & $16$ & $32$ & $64$ & $128$ & $256$ \\
\midrule
$\tau$ basis        & $45.9$ & $49.9$ & $52.8$ & $55.7$ & $57.6$ & $59.0$ \\
cross-term basis    & $45.3$ & $48.8$ & $51.4$ & $53.5$ & $54.8$ & $55.7$ \\
\bottomrule
\end{tabular}
\end{table}

The separate cross-term-only heuristic uses
$\sum_{t\ne s}\Delta W_t^\top\Delta W_s=\tau^\top\tau-
\sum_t\Delta W_t^\top\Delta W_t$. This does not improve compressed-merge accuracy: the
$\tau$ basis wins at every tested budget, the gap widens with
$r$, and at $r=256$ it wins on all eight tasks separately. Concentrating empirical cross-task
coupling is therefore insufficient for better compressed-merge accuracy in this experiment.

Calibration is close to rank-preserving. Across the $190$ ViT-B/32 pairs, raw and calibrated
overlap have Spearman $0.995$ and share $17$ of the top $20$ pairs; ViT-B/16 gives $0.994$ and
shares $19$ of the top $20$. Thus an ordering-only consumer receives nearly the same ordering on these pools,
although we do not retrain the mergeability predictor of \citet{demystify2026}.

\section{Estimating $C_H$}
\label{app:cov}

$C_H$ is $d_{\mathrm{in}} \times d_{\mathrm{in}}$, which is $1.4$ GB per module at
$d_{\mathrm{in}} = 18944$ and infeasible to store across a network.
For the vision encoders studied here the exact matrix is affordable ($594$ MB across all $72$
modules of ViT-B/32) and we use it.
For larger models we provide a two-pass estimator: a Frequent Directions sketch locates the
subspace, and a second pass accumulates the exact $m_C \times m_C$ second moment inside it.
The second pass is necessary. Frequent Directions biases eigenvalues
downward by construction, and whitening applies $C_H^{-1/2}$, so an understated eigenvalue is an
over-amplified direction; on synthetic data at sketch size $128$ the eigenvalue error falls from
$0.49$ to $0.0016$ between the two passes while the top-$8$ subspace agreement is $0.998$
either way.

\paragraph{The calibration set is part of the measurement.}
$C_H$ is an expectation over some input distribution, and the answer depends on which.
We draw it from a mixture across several of the benchmark's own datasets, since a single one
would make $C_H$ the base geometry \emph{on that task's inputs} and
reintroduce the conditioning the test exists to remove.
A separate rebuilding record lists all image identifiers, per-source counts, hashes, and fold
memberships, while each covariance manifest records its item count and a basename--byte-size
hash. We verified that the two records enumerate the same $1001$ basenames; the release links
them explicitly rather than implying that source labels are embedded in each flattened
covariance manifest. The covariances themselves come to $26$ GB and need not be distributed; a
reproducibility artifact must include this rebuilding record and the extraction commands.

\section{Reproducibility record}
\label{app:repro}

\paragraph{Checkpoints and task pools.}
The frozen bases are \texttt{openai/clip-vit-base-patch16},
\texttt{openai/clip-vit-base-patch32}, and \texttt{openai/clip-vit-large-patch14}. Full-fine-tuned
specialists follow \texttt{tanganke/<architecture>\_<task>}. The $20$ task labels are
\texttt{cifar10}, \texttt{cifar100}, \texttt{dtd}, \texttt{emnist\_letters}, \texttt{eurosat},
\texttt{fashion\_mnist}, \texttt{fer2013}, \texttt{food101}, \texttt{gtsrb}, \texttt{kmnist},
\texttt{mnist}, \texttt{oxford-iiit-pet}, \texttt{oxford\_flowers102}, \texttt{pcam},
\texttt{rendered-sst2}, \texttt{resisc45}, \texttt{stanford-cars}, \texttt{stl10},
\texttt{sun397}, and \texttt{svhn}. The matched eight-task comparison uses \texttt{dtd},
\texttt{eurosat}, \texttt{gtsrb}, \texttt{mnist}, \texttt{resisc45}, \texttt{stanford-cars},
\texttt{sun397}, and \texttt{svhn}; its adapter repositories append \texttt{\_lora-16} or
\texttt{\_l-lora-16}. Rank and scale are read from each \texttt{adapter\_config.json}, and
untargeted modules are dropped rather than treated as zero subspaces.

\paragraph{Calibration and numerical settings.}
The rebuilding record contains $143$ images from each of \texttt{cifar10}, \texttt{dtd},
\texttt{eurosat}, \texttt{gtsrb}, \texttt{resisc45}, \texttt{stl10}, and \texttt{sun397}; fold
sizes are $126,125,125,125,125,125,125,125$. Images are the first valid examples encountered in
the streaming training split of each \texttt{tanganke/<source>} dataset, converted to RGB.
Extraction uses each model's \texttt{CLIPImageProcessor}, all sequence tokens (class and patch), exact uncentred
$\sum h h^\top/n$ accumulation in float32, $m_C=64$, and CPU execution. The extractor-input
basename/size hash common to the three covariance manifests has prefix \texttt{8f88002f7557}; the separate logical image/fold
record has identifier-list-hash prefix \texttt{73797f525986}. Full digests are in the artifact.

\paragraph{Randomisation and verification.}
The activation-null scalar results are the exact expectation in Eq.~\ref{eq:null-expectation};
they use no random draws. On the matched full-fine-tuning pool, a sampled check gives null
$0.02330$ versus exact $0.02333$. External-reference diagnostics use finite draws, the two
factor-null adapter arms use four independent $16$-draw batches (seeds $0$, $1000003$,
$2000003$, and $3000017$), pooled row-wise to $64$ draws, and projector rebuilds use six.
The batch-based pooled Monte Carlo standard errors (sample SD divided by $\sqrt 4$, three degrees
of freedom) are $0.0067$ and $0.0054$ percentage points for LoRA-16 and linearized LoRA-16;
pooling changes their seed-$0$ $16$-draw fractions by $-0.0109$ and $-0.0093$ points.
The common-reference baseline uses $20$ pairs for each of eight fixed modules at seed $0$.
\texttt{analysis/check\_paper\_numbers.py} recomputes point estimates from row-level CSVs, and
\texttt{analysis/overlap\_bootstrap.py} reproduces the $50{,}000$ crossed-resampling stability intervals.
The artifact pins immutable revisions for all bases, specialists, adapters, and calibration
datasets, together with a locked runtime environment. The original B/16 covariance manifest
omitted its base-revision label; its two candidate revisions are tensor-identical in all $400$
tensors, with content hashes recorded, so the numerical weights are unambiguous. Large model
weights, covariance tensors, and calibration images are not bundled.

\section{How the three defects were fixed}
\label{app:defects}

\paragraph{The complement was rotated inside a fixed frame.}
$\Delta W$ splits into a part inside $\mathrm{span}(V)$ and a residual $R$ whose rows lie in the
complement. The first version drew a random $128$-dimensional frame of that complement and
rotated $R$ inside it, leaving the remainder bit-identical between two task vectors. The
complement is $d_{\mathrm{in}} - m_C$ dimensional, so this covered $18\%$ of it at
$d_{\mathrm{in}}=768$ and $4\%$ at $d_{\mathrm{in}}=3072$, while about three quarters of the real
read mass lives there. The fix needs no square Haar matrix: $R$ has rank at most
$\min(d_{\mathrm{out}}, d_{\mathrm{in}}-m_C)$, so a Haar map on the complement carries its right
singular vectors to a uniform frame of that rank, which one thin SVD and one QR deliver exactly.

\paragraph{Singleton blocks were skipped.}
Blocks of size one were left alone on the grounds that a one-dimensional space cannot be rotated.
Its orthogonal group is $\{+1,-1\}$, whose Haar measure is a fair coin. Independent flips preserve
diagonal coordinate masses and randomise off-diagonal cross-coordinate signs, which cancel in
expectation, while leaving $\Delta W\Delta W^\top$ and hence the spectrum untouched. On ViT-B/16 the singletons
carry $71\%$ of the stored eigenvalue mass, so skipping them retained substantial coupling.

\paragraph{The LoRA generator matched a norm where a factor spectrum was required.}
The factor-conditioned null of \S\ref{sec:lora-null} preserves $A$ and the singular values of
$B$; it does not generally preserve the product spectrum of $BA$. Its first implementation
redrew $B$ as an i.i.d. Gaussian at matched Frobenius norm, leaving even the factor spectrum
uncontrolled: a Gaussian is nearly Marchenko--Pastur where trained $B$ decays. On fitted-reference
pools with shared $A$ and independent $B$, the flat-spectrum generator reports $+0.33$ when $B$
has no decay and $-0.16$ when it does. Drawing $B$ at its observed singular values holds the
reference offset within $+0.008$ across the same sweep. This removes the observed factor-spectrum
mismatch; it is not a claim of exact product-spectrum preservation.
Indeed the nonzero squared singular values of the redraw are the eigenvalues of
$\operatorname{diag}(S_B)W^\top(AA^\top)W\operatorname{diag}(S_B)$, which depend on $W$ unless
$AA^\top$ is isotropic. A uniform product-spectrum-preserving rotation inside
$\mathrm{row}(A)$ fails differently: it discards this anisotropy and has constructed-reference
offset $+0.2023$. Rotating only to $BQ$ retains a learned singular
frame and changes the product singular values by $19.7\%$.

\section{Projector rebuild}
\label{app:projector}

\begin{table}[h]
\centering
\caption{Projector-basis overlap after independent draws from each fitted conditional null, and
the observed basis's mass in the top-$64$ activation subspace. ViT-B/16, eight tasks, $r_p=128$
over $12$ full-fine-tuning modules and $r_p=16$ over $24$ adapter modules.}
\label{tab:projector}
\begin{tabular}{@{}lrrrr@{}}
\toprule
pool & observed--null & chance & in activation & chance \\
\midrule
full fine-tuning & $0.2594$ & $0.1667$ & $0.2612$ & $0.0833$ \\
LoRA-16          & $0.8750$ & $0.0208$ & $0.1356$ & $0.0833$ \\
\bottomrule
\end{tabular}
\end{table}

\citet{qiu2026nufilt} filter an incoming task vector with $P=I-VV^\top$, where $V$ contains the
top-$r_p$ right singular vectors of the cumulative update $\tau$. We build this basis once from
the observations and once after independently applying the regime-specific null to every task.
The activation null preserves full-update spectra and occupancies; the LoRA null preserves its
factor-level invariants. The second column reports overlap of the resulting bases, while the
fourth is a separate concentration statistic. Neither identifies how that structure arose, tests
accuracy, or contradicts the gains reported for the projector. A basis reproducible
under a conditional null cannot, without further evidence, identify the cross-task mechanism used
to justify it. The LoRA value has no projector-level positive control; the factor-null overlap
stress test covers only one different, narrow intervention.

\section{Which correction did the work}
\label{app:ablation}

The two corrections to the activation null can be switched independently. Scoring the same
planted pairs at $s=2$ with each combination separates them.

\begin{table}[h]
\centering
\begin{tabular}{@{}llr@{}}
\toprule
full complement & singleton flip & level-centred recovery \\
\midrule
        &        & $49.0\%$ \\
        & \ding{51} & $77.7\%$ \\
\ding{51} &        & $89.2\%$ \\
\ding{51} & \ding{51} & $99.3\%$ \\
\bottomrule
\end{tabular}
\end{table}

Each arm is centred by its own $s=0$ level and divided by the same delivered increment as in
Eq.~\ref{eq:recovery}. The singleton flip alone raises recovery from $49.0\%$ to $77.7\%$, the
full complement alone to $89.2\%$, and both to $99.3\%$; the effects overlap.
The singleton flip is not a token gesture. On ViT-B/16 the singletons hold $71\%$ of the stored
eigenvalue mass, and leaving them fixed exempted most of $\mathrm{span}(V)$.

\section{Planted-structure recovery, by level}
\label{app:power}

\paragraph{Planting operator.}
Let $P_0=U_0\operatorname{diag}(\sigma_1,\ldots,\sigma_\rho)V_0^\top$ be a thin SVD with
descending singular values, and let $S=V_A[:,1{:}s]$ contain the source task's first $s$ right
singular vectors. We set the first $s$ columns of $\widehat V$ to $S$. In order, we project
$V_0[:,s{+}1{:}\rho]$ onto $S^\perp$ and apply QR/Gram--Schmidt against all columns already
accepted; any algebraic or numerical rank deficiency is filled by a Haar frame in the remaining
complement. We then set
$P_s=U_0\operatorname{diag}(\sigma_1,\ldots,\sigma_\rho)\widehat V^\top$. Hence
$\widehat V^\top\widehat V=I$ and every singular value and left singular vector of $P_0$ is
retained. Away from a tie at the relevant singular-value boundary, the first $s$ leading right
directions match the source exactly. The same planted matrices are scored by every candidate null.

\begin{table}[h]
\centering
\caption{Positive-control response on ViT-B/16, four modules, six tasks. \texttt{initial} is
the first activation null and \texttt{corrected} is \S\ref{sec:null}. Both score the same planted
pairs. \emph{delivered} is
$(\texttt{raw}(s)-\texttt{raw}(0))/[(s/k)(1-\mathcal O_0)]$. Recovery is the
level-centred excess increment divided by the delivered overlap increment, Eq.~\ref{eq:recovery};
it is therefore comparable across planting levels and is not the uncentred excess divided by
$s/k$.}
\label{tab:power}
\begin{tabular}{@{}lrrrr@{}}
\toprule
planted & \texttt{raw} & delivered & initial recovery & corrected recovery \\
\midrule
$0/8$ & $0.02771$ & --- & $X_0=+0.00004$ & $X_0=-0.000623$ \\
$1/8$ & $0.14129$ & $93.45\%$ & $45.6\%$ & $99.6\%$ \\
$2/8$ & $0.26235$ & $96.53\%$ & $49.5\%$ & $99.4\%$ \\
$4/8$ & $0.50557$ & $98.30\%$ & $51.7\%$ & $99.3\%$ \\
\bottomrule
\end{tabular}
\end{table}

\section{Overlap decomposition, in full}
\label{app:main}

\begin{table}[h]
\centering
\caption{Cross-task read-subspace overlap and its scalar decomposition. Complete encoder pools
use all $190$ task pairs; the matched eight-task row uses all $28$ pairs. \texttt{raw} is
Eq.~\ref{eq:overlap}, \texttt{iso} is $k/d$, and \texttt{null} is the exact expectation in
Eq.~\ref{eq:null-expectation}. Terms are averaged before computing the reproduced fraction. No
Monte Carlo draw enters these scalar results. ViT-L/14 is run on \texttt{q} and \texttt{v} only.}
\label{tab:main}
\small
\setlength{\tabcolsep}{4pt}
\begin{tabular}{@{}lrrrrrr@{}}
\toprule
& modules & \texttt{raw} & \texttt{iso} & \texttt{null} & residual & reproduced \\
\midrule
ViT-B/16 & $72$ & $0.104574$ & $0.009115$ & $0.033309$ & $0.071265$ & $25.3451\%$ \\
ViT-B/32 & $72$ & $0.095313$ & $0.009115$ & $0.028420$ & $0.066893$ & $22.3961\%$ \\
ViT-L/14 $q/v$ & $48$ & $0.071402$ & $0.007813$ & $0.015377$ & $0.056024$ & $11.8961\%$ \\
matched FT8 $q/v$ & $24$ & $0.089930$ & $0.010417$ & $0.023333$ & $0.066597$ & $16.2440\%$ \\
\bottomrule
\end{tabular}
\end{table}

For the exploratory ViT-B/16 role grouping in Figure~\ref{fig:main}, the same unadjusted ratio is
$16.5568\%$ for residual-reading modules (\texttt{q}, \texttt{k}, \texttt{v}, \texttt{fc1}) and
$38.9609\%$ for block-internal readers (\texttt{out\_proj}, \texttt{fc2}).

On the matched FT8 pool, a sampled-matrix validation gives mean null $0.02330$, versus the exact
$0.02333$. Sampled matrices remain necessary for projector and positive-control experiments.

\begin{table}[h]
\centering
\caption{External-reference \texttt{H0} stress diagnostics. These activation-generated references
are not the exact Haar self-orbit, whose population \texttt{H0}$=0$ by identity. The finite-draw
offsets can mix reference mismatch and over-destruction and are never added to an estimate.}
\begin{tabular}{@{}lr@{}}
\toprule
pool and candidate action & realised external \texttt{H0} \\
\midrule
matched FT8, activation null & $-0.00040$ \\
LoRA-16, activation null & $-0.00037$ \\
lin-LoRA, activation null & $-0.00037$ \\
\bottomrule
\end{tabular}
\end{table}

\section{The generative null and why it fails}
\label{app:generative}

The obvious construction is generative: draw a random update whose input-side covariance is
proportional to the base activation covariance, adding an isotropic draw in the unresolved
complement. It does not work because it must invent the relationship between the trained update
and $C_H$. Independent draws then concentrate on the same high-variance activation directions
even when observed updates have a different occupancy profile. On the explicitly scoped
eight-task, eight-module common-reference diagnostic in Table~\ref{tab:baselines}, this
complement-aware generator reports native raw-minus-null excess $-0.8959$, whereas the
observed-profile-preserving null reports $-0.0007$. These are baseline-specific stress levels,
not estimates of the headline conditional fraction. An earlier exploratory generator omitted
the unresolved complement entirely; we discard its quantitative output rather than using a
known-defective construction. The comparison locates the requirement precisely: a conditional
null here must retain the observed update--activation occupancy rather than generate that
relationship from an unverified model.
